\section{Introduction}

Precision Livestock Farming (PLF) increasingly relies on computer vision (CV) systems to monitor and manage livestock. Localizing individual animals from images and videos helps farm operators optimize health, behavior monitoring, and resource management. However, despite progress in applying computer vision to livestock settings, a critical challenge persists in deploying models across diverse environments. Inconsistent lighting, varied camera angles, and differences in farm layouts can significantly affect model generalization.

This study focuses on the generalization capabilities of YOLO-family object detection models in the task of indoor cow localization. Specifically, we investigate how environmental factors and model complexities impact performance and explore whether fine-tuning with domain-specific data improves detection outcomes. To support this study, we introduce the \textbf{COws LOcalization (COLO)} dataset, a public resource for advancing cow detection in complex barn environments.

We define three central hypotheses to guide this investigation:

\begin{itemize}
    \item \textbf{Model generalization is equally influenced by changes in lighting conditions and camera angles.} Should camera angles prove more impactful than lighting conditions, it would be advisable to prioritize camera placement when deploying CV models in new environments.
    \item \textbf{Higher model complexity guarantees better generalization performance.} If a highly complex model does not ensure superior performance, future studies might consider adopting less computationally demanding models that still enhance performance.
    \item \textbf{The advantages of using fine-tuned models as initial training weights are persistent over pre-trained models.} If the advantages diminish as the training sample size increases in a similar cow localization task but different environments, the fine-tuning efforts may be deemed unnecessary when the deployment environment varies over multiple locations on a farm.
\end{itemize}

To facilitate these investigations, a public dataset named COws LOcalization (COLO) \cite{COLODataset2023} will be developed and made available to the community. The findings of this study are expected to provide practical guidelines for Precision Livestock Farming (PLF) researchers on deploying CV models, considering available resources and anticipated performance.

\section{Background Study}

\subsection*{Object Localization and Its Applications}

Localizing livestock individuals from images or videos has become an essential task in precision livestock farming (PLF) \cite{fernandes2020image}. Such techniques allow farm operators to manage animal well-being and health in real-time, optimizing their resource management and improving sustainability \cite{morrone2022industry, hao2023cattle}. Technically speaking, in the field of computer vision (CV), which is a subfield of artificial intelligence (AI) that focuses on translating visual information into actionable insights, localization tasks can be further categorized into object detection, object segmentation, and pose estimation. Object detection is the simplest form among these tasks, localizing objects of interest by enclosing them within a rectangular bounding box defined by x and y coordinates, pixel width, and pixel height \cite{viola2001rapid}. Successful instances in this category include YOLO (You Only Look Once) \cite{redmon2016you}, Faster R-CNN (Region Convolutional Neural Networks) \cite{girshick2015fast}, and SSD (Single Shot MultiBox Detector) \cite{liu2016ssd}. These models have been adopted and applied by animal scientists for detection in precision livestock farming. For example, a study \cite{yu2022automatic} leveraged the DRN-YOLO model \cite{xu2020improved} to predict the eating behavior of dairy cows. This approach automates the assessment of feeding behavior, a critical indicator of cow health and productivity, and has saved labor efforts in complex farm settings. Another notable work is presented in \cite{nasirahmadi2019deep}, where the authors developed a posture detection system for pigs using deep learning models such as Faster R-CNN, SSD, and R-FCN, coupled with 2D imaging. This system accurately identifies standing and lying postures of pigs under commercial farm settings.

To achieve finer localization, object segmentation is employed to outline object contours pixel-wise, while pose estimation is performed by orienting and marking the key points of the object \cite{hariharan2015hypercolumns}. Some popular object segmentation models include Mask R-CNN \cite{he2017mask}, MS R-CNN \cite{huang2019mask}, and U-Net \cite{siddique2021u}. This method of segmentation has also been applied in the field of PLF. In the study \cite{noe2022automatic}, the authors developed a method using Mask R-CNN \cite{he2017mask} to segment and outline cattle in feedlots. Their technique enhances images and extracts key frames to accurately detect cattle, achieving superior precision with a mean pixel accuracy of 0.92. This supports advanced, real-time monitoring of cattle in PLF. Another study group \cite{tu2021automatic} developed the PigMS R-CNN framework \cite{huang2019mask} to enhance the monitoring of group-housed pigs. This framework employs a 101-layer residual network along with a feature pyramid network and soft non-maximum suppression to effectively detect and segment pigs, thereby improving the accuracy of identifying and locating individual pigs in complex environments.

\subsection*{Model Generalization, Pre-Training, and Fine-Tuning}

Although implementing image-based systems in livestock production has become more common, current studies primarily focus on accuracy in homogenous environments and rarely address the challenges of model generalization. How a model can generalize to new environments is critical when farm operators deploy existing CV models in their own settings. Good generalization performance ensures that the model can reproduce similar results as reported in the original study, even in new environments with different conditions. Factors such as camera angles and the presence of occlusions can impact generalization in the deployment environment. Deploying the same model in a new environment does not necessarily guarantee the same performance as reported in the original study. Li et al. \cite{li2021practices} also pointed out that the lighting conditions on farms in real applications can be highly variable, leading to poor generalization performance.

One explanation for poor generalization is the discrepancy between the pre-training process and the specific use case. Most CV models are released with pre-trained weights, obtained from training on a large-scale dataset. For example, the COCO dataset \cite{lin2014microsoft} is a general-purpose dataset containing over 200,000 images and a wide range of object categories, such as vehicles and household items. Directly deploying a model pre-trained on the COCO dataset to detect cows in a farm setting may not ensure satisfactory performance, as the dataset does not contain enough cow instances in different view angles or occlusions. To alleviate this discrepancy, fine-tuning is a common practice that modifies the prediction head of the pre-trained model and updates the weights on a new dataset more relevant to the specific use case. Most application studies have adopted this approach to improve model generalization on their specific tasks \cite{han2021pre,guirguis2022cfa,gupta2023novel}.

Nevertheless, fine-tuning is not guaranteed to be successful, as the outcome depends on both the quantity and quality of the annotated dataset. For example, Zin et al. \cite{zin_automatic_2020} deployed an object detection model to recognize cow ear tags in a dairy farm. Although the model achieved a high accuracy of 92.5\% in recognizing the digits on the ear tags, more than 10,000 images were required for fine-tuning. Assembling such a large dataset is labor-intensive and requires specific training in annotating the images. The annotated dataset is rigorously organized in a specific format, such as COCO or YOLO, each with its own structure and coordinate systems, requiring professional annotation tools like Labelme \cite{labelme2023}, CVAT \cite{cvat2023}, or Roboflow \cite{roboflow2023}.

\subsection*{YOLO Models}

Before YOLO, object detection methods typically involved either using ``sliding windows with classifier'' or ``region proposals with classifier.'' The sliding windows method required running the classifier hundreds or thousands of times per image. On the other hand, advanced region proposal-based approaches divided the task into two steps: first, identifying potential object regions (i.e., region proposals) and then applying a classifier to these regions. In contrast, YOLO models are capable of performing object detection in a single pass through the network, which is why the acronym YOLO stands for ``You Only Look Once.''

% The sliding window approach involved scanning the image at multiple scales and locations using a fixed-size window, and then applying a classifier (e.g., SVM or boosted decision trees) to each window to determine whether it contained an object. This method, while straightforward, was computationally expensive due to the exhaustive number of windows that had to be evaluated. Notable early work using this method includes the Viola-Jones face detector (Viola & Jones, 2001), which pioneered real-time object detection using Haar-like features and a cascade of classifiers.

% Later, region proposal-based approaches improved efficiency and accuracy by decoupling object localization from classification. These methods first generated a small set of region proposals—likely object-containing areas—using algorithms like Selective Search (Uijlings et al., 2013), and then applied a convolutional neural network (CNN) to classify each proposed region. A landmark example is R-CNN (Girshick et al., 2014), followed by its faster successors Fast R-CNN (Girshick, 2015) and Faster R-CNN (Ren et al., 2015), the latter introducing Region Proposal Networks (RPNs) to make the proposal process learnable and more efficient.

% In contrast to these multi-stage pipelines, YOLO (You Only Look Once) (Redmon et al., 2016) revolutionized object detection by reframing it as a single regression problem. Instead of generating region proposals, YOLO divides the image into a grid and predicts bounding boxes and class probabilities directly from the full image in a single forward pass of the network. This design enables real-time detection with a significant speed advantage, albeit initially at the cost of some localization accuracy. Subsequent versions, including YOLOv2 (Redmon & Farhadi, 2017), YOLOv3 (Redmon & Farhadi, 2018), and more recently YOLOv4, YOLOv5, and YOLOv8, have made strides in balancing speed and accuracy for diverse applications.
% References:
% 	•	Viola, P., & Jones, M. (2001). Rapid object detection using a boosted cascade of simple features. CVPR.
% 	•	Uijlings, J. R. R., et al. (2013). Selective search for object recognition. IJCV.
% 	•	Girshick, R. (2014). Rich feature hierarchies for accurate object detection and semantic segmentation. CVPR (R-CNN).
% 	•	Girshick, R. (2015). Fast R-CNN. ICCV.
% 	•	Ren, S., et al. (2015). Faster R-CNN: Towards real-time object detection with region proposal networks. NeurIPS.
% 	•	Redmon, J., et al. (2016). You Only Look Once: Unified, real-time object detection. CVPR.
% 	•	Redmon, J., & Farhadi, A. (2017). YOLO9000: Better, faster, stronger. CVPR.
% 	•	Redmon, J., & Farhadi, A. (2018). YOLOv3: An incremental improvement. arXiv:1804.02767.


%\textbf{YOLOv11} introduced significant architectural enhancements to improve performance and efficiency. Key innovations include the introduction of the C3K2 (Cross Stage Partial with kernel size 2) block, Spatial Pyramid Pooling - Fast (SPPF), and Convolutional block with Parallel Spatial Attention (C2PSA). These components enhance feature extraction and computational efficiency, resulting in improved mean Average Precision (mAP) and versatility across various computer vision tasks, including object detection, instance segmentation, pose estimation, and oriented object detection. Additionally, YOLOv11 offers multiple model sizes, from nano to extra-large, catering to diverse application needs from edge devices to high-performance computing environments\cite{ultralytics2024yolo11, khanam2024yolov11}.

%\textbf{YOLOv12} represents a shift towards an attention-centric architecture, departing from traditional convolutional neural network (CNN)-based approaches. It incorporates an optimized backbone (R-ELAN), 7×7 separable convolutions, and FlashAttention-driven\cite{dao2023flashattention, d} area-based attention mechanisms. These innovations enhance feature extraction, efficiency, and detection robustness. YOLOv12 achieves higher mAP and faster inference speeds compared to its predecessors, making it suitable for applications requiring real-time performance and high accuracy\cite{ultralytics2024yolo12, tian2025yolov12}.

%\textbf{DETR-RT} (Detection Transformer - Real-Time) is a transformer-based object detection model designed for real-time applications. It eliminates the need for traditional components like non-maximum suppression by formulating object detection as a direct set prediction problem. DETR-RT leverages transformer architectures to model global relationships within an image, enabling robust detection of objects with varying scales and in complex scenarios. Its end-to-end approach simplifies the detection pipeline and improves adaptability to dynamic environments\cite{ultralytics2024detrt, carion2020end}.

%%%%%%%%%%%%%%%%% In my words %%%%%%%%%%%%%%%%%%%%%%%


\textbf{YOLOv11} introduces several architectural improvements to enhance both speed and accuracy. Key innovations include the C3K2 block (Cross Stage Partial with kernel size 2) for better feature extraction, Spatial Pyramid Pooling - Fast (SPPF) to improve multi-scale detection, and the Convolutional block with Parallel Spatial Attention (C2PSA) to help the model focus on important image regions. These enhancements make YOLOv11 more efficient and effective for tasks such as object detection, instance segmentation, pose estimation, and oriented object detection. Additionally, YOLOv11 is available in multiple sizes, ranging from lightweight models for edge devices to larger versions for high-performance applications \cite{ultralytics2024yolo11, khanam2024yolov11}.

\textbf{YOLOv12} takes a different approach by incorporating attention mechanisms to improve object detection. Instead of relying only on traditional Convolutional Neural Networks (CNNs), it uses an optimized backbone (R-ELAN), 7×7 separable convolutions, and FlashAttention-driven area-based attention mechanisms \cite{dao2023flashattention, dao2022flashattention}. These features allow YOLOv12 to process images more efficiently, capture important details better, and perform well even in complex scenes. As a result, YOLOv12 achieves higher accuracy and faster inference speeds than its predecessors, making it ideal for real-time applications \cite{ultralytics2024yolo12, tian2025yolov12}.

\textbf{DETR-RT} (Detection Transformer - Real-Time) is a transformer-based object detection model designed for speed and accuracy. Unlike traditional YOLO models, it does not require predefined anchor boxes or post-processing techniques like non-maximum suppression. Instead, DETR-RT directly predicts object locations and classes using a transformer architecture. This allows it to understand object relationships in an image better, making it more robust for detecting objects of varying sizes and in complex environments \cite{ultralytics2024detrt, carion2020end}.

In our study, we combine both CNN-based models (YOLOv11) and transformer-based architectures (YOLOv12 and DETR-RT) to balance speed, accuracy, and efficiency for real-time object detection.

\subsection*{Public Datasets}

A public dataset helps the community to develop methodology based on the same baseline. One famous example in computer vision is the ImageNet dataset \cite{deng2009imagenet}, which serves as a benchmark for image classification. AlexNet \cite{krizhevsky2017imagenet}, the winner of the ImageNet Large Scale Visual Recognition Challenge in 2012, demonstrated its outstanding capability to classify images in the ImageNet dataset using Rectified Linear Units (ReLU) as the activation function, rather than the traditional sigmoid function. The success of AlexNet accelerated the development of CV models in the following years, such as VGG \cite{karen2014very}, GoogLeNet \cite{szegedy2015going}, ResNet \cite{targ2016resnet}, and DenseNet \cite{huang2017densely}. However, similar to the challenges that pre-trained models face in specific use cases, a generic public dataset, such as ImageNet \cite{deng2009imagenet} and COCO \cite{lin2014microsoft}, may not be sufficient for PLF applications.

There have been efforts to create public datasets for livestock scenarios. For example, the CattleEyeView dataset was collected to support applications like cattle pose estimation and behavior analysis, providing extensive annotations across 30,703 frames from top-down video sequences of cows \cite{ong2023cattleeyeview}. Another study \cite{t2020long} leverages a public dataset for pigs comprising 3600 images from 12 videos of group-housed pigs. The dataset is particularly designed for applications such as pig tracking. Additionally, the ``OpenCows2020'' dataset, developed by researchers from the University of Bristol, is a public dataset specifically designed for advancing non-intrusive monitoring of cattle. It supports precision farming applications such as automated productivity assessment, health and welfare monitoring, and veterinary research, including behavioral analysis and disease outbreak tracing. The dataset consists of 11,779 images with 13,026 labeled objects, mainly focusing on cattle \cite{visualization-tools-for-opencows2020-dataset}.